\chapter{Sanntidsanalyse}
\label{ch:timing}
\section{Modellen må svare til programmet}
Den aktive kjernen er sekvensiell og inneholder venting på tid og I/O.
Én generell periodisk CPU-modell kan derfor ikke uten videre bevise
reaksjonen. Den ønskede lysperioden er 50 ms med lav prioritet; den er et
mykt mål og kan få jitter fra høyere prioriterte aktiviteter. Interne
driveravbrudd, taktkilde, prioritet og blokkering må dokumenteres ved måling.

For fastprioriterte periodiske/sporadiske oppgaver brukes ofte iterasjonen
\begin{equation}\label{eq:rta}
 R_i^{(k+1)}=C_i+B_i+
 \sum_{j\in hp(i)}\left\lceil\frac{R_i^{(k)}+J_j}{T_j}\right\rceil C_j,
\end{equation}
med WCET $C$, blokkering $B$, periode/minste ankomstavstand $T$ og jitter
$J$ \cite{josephpandya,rta}. Modellen krever begrunnede forutsetninger;
selvsuspensjon og I/O-venting må behandles særskilt. Kriteriet er
$R_i\le D_i$. Her foreligger verken dokumenterte $C_i$ eller et komplett
tasksett som lar oss hevde at dette er oppfylt.

\section{Konfigurerte tidsbidrag}
\begin{tabularx}{\textwidth}{@{}P{42mm}P{29mm}Y@{}}\toprule
Bidrag & Kodeparameter & Tolkning\\\midrule
Akustisk triggeravstand & Minst 65 ms & Nominelt frontbesøk hver 130 ms, side hver 260 ms i Cruise. Forsinkelser kan øke intervallene.\\
ECHO-frister, applikasjon & 3 + 23 ms & Grense for stigning og fall i aktiv driver; GPIO-diagnostikk bruker 10 + 23 ms.\\
I2C-skriving & 20 ms per kall & Timeoutparameter, ikke alene bevist kjøretid.\\
Fallback-stopp & Opptil 8 kall & 160 ms samlet konfigurasjonsbudsjett, pluss overhead og scheduling.\\
Servo stabilisering & 500 ms & Stasjonær designventing, ikke målt maksimal servotid.\\
Sideevidens & Høyst 2000 ms & Alder målt fra eldste stråle frem til svingkommando.\\
Lys-task & 50 ms & Ønsket myk oppdatering, ingen målt jittergrense.\\\bottomrule
\end{tabularx}

Ved fremoverkjøring må en hindring først vente til neste relevante trigger.
Deretter kommer sensorlesing, beslutning og stopptransaksjon:
\begin{equation}
 t_r=G_{\mathrm{relevant}}+L_{\mathrm{sensor}}+
 C_{\mathrm{beslutning}}+L_{\mathrm{stopp}}+I,
\end{equation}
der $G$ er største målte/analytisk begrunnede triggergap og $I$ omfatter
avbrudd, scheduling og andre forsinkelser. 130 ms er et nominelt
frontintervall, ikke automatisk en øvre grense for $G$.

\section{Stoppknapp og fysiske deadlines}
Knapp B er en \emph{programvarebasert stoppforespørsel}. Avbruddet setter
flagget raskt, men driverkall og en pågående venting kan forsinke
foreground-kodens reaksjon. I Cruise sjekkes flagget før og etter
sensorlesing samt i motorpakken før og etter sending. Under stasjonært
servoarbeid kan programmet fortsatt være i en 500 ms venting.
Knappen er derfor ikke dokumentert som en uavhengig hard nødstopp.
Ved bortfall av I2C kan programmet ikke garantere å slå av motorene.

En fysisk reaksjonsfrist må velges fra tilgjengelig klaring og
ligning~\ref{eq:stop}, med en konservativt målt fart og retardasjon.
\todo{Fyll inn $v_{\max}$, minste dokumenterte $a_b$, margin $m$, største
reaksjonstid og valgt deadline med rådata.} En fast terskel på 45 cm kan
ikke forsvares bare fordi den finnes i kildekoden. Svinggrensen 25 cm og
sidegrensen 15 cm må tilsvarende vurderes mot chassis og måleusikkerhet.

\section{Plan for tidsmåling og verifikasjon}
Mål på micro:biten med tidsstempel eller logikkanalysator: hindrings-/
stoppstimulus, relevant trigger, ferdig ECHO, stoppforespørsel,
akseptert motorstopp og faktisk stillstand. Skill prosessortid fra
veggklokketid og fysisk bremsing. Registrer også triggergap, lysjitter,
skannealder ved kjøring og stakkmargin. Logging som endrer lasten må
beskrives i forsøksoppsettet.

Gjenta ved begge svingretninger, laveste forventede batterispenning og
relevante reflektorer. Rapporter antall forsøk, maksimum, fordeling og
usikkerhet. Det største observerte tallet er et måleresultat, ikke bevist
WCET. Hvis fristen overskrides, må fart eller driftsområde reduseres,
eller arkitekturen endres før kravet godkjennes.
